\scanpage{223}

\mspage{97}

base, et c'est assuré par 7.2.3. c) (dont on avait déjà besoin dans 7.1., 3°, pour établir $C\cap W\subset\deleted{\unsure{C_{\mathrm{stab}}}}\widetilde{\mathrm{TC}_0}$. (\uline{NB} les hypothèses de 7.2. sont pratiquement celles de 7.1, 3°, sauf qu'on suppose que $W$ est un localiseur (\unsure{faiblement|au sens} : modérément saturé), de façon précise on n'a utilisé, jusqu'à présent, comme hyp. suppl., que la stabilité de $W$ par composition.)

\deleted{Si on rajoute les hypothèses \unsure{a) à|a) et}}

Il reste à établir que $(W,C=\mathrm{Mon.},F=(\mathrm{TC})_*)$ est un triple de Quillen clos, ssi les conditions 7.2.6 et (7.2.7), i.e.

\[
\left\{\begin{aligned}
&C\inserted{=\mathrm{Mon}.}\ \text{stable par cochangement de base}\\
&\mathrm{TF}=C_*=(\mathrm{Mon})_*\subset W
\end{aligned}\right.\tag*{(7.2.7)}
\]

sont satisfaites. La nécessité est évidente, par définition des triples de Quillen clos, et pour la suffisance, on note qu'il résulte du th. 7.1, \unsure{2°|3°} et de \deleted{la} partie déjà prouvée de 7.2, que $(C,\mathrm{TF}),(\mathrm{TC},F)$ sont des couples de Quillen clos, et on a de plus

\[
\mathrm{TC}=C\cap W,\quad \mathrm{TF}\subset F\cap W
\]

(puisque $\mathrm{TF}\subset W$ par 7.1.7, et $\mathrm{TF}\subset\unsure{F|R}$ \deleted{\unsure{puisque}} $\mathrm{TC}\subset C$). \deleted{Il s'ensuit} On sait
% archive sheet number: 222
